\section{Antes de empezar}
\subsection*{Cómo entras}
\begin{enumerate}
  \item \emph{Titulación} → \emph{Validar trámites} (bandeja) y, al final, \emph{Expedientes a Servicios Escolares}.
  \item La Jefatura firma su visto bueno en \emph{Firmas pendientes}.
\end{enumerate}
\begin{tencuenta}
\begin{itemize}
  \item Al devolver un documento la observación es obligatoria; el pasante la ve arriba de su trámite.
  \item Los documentos generados por la plataforma ya vienen firmados; revísalos y valídalos igual que los demás.
  \item Cada renglón de la verificación final se marca como entregado o \emph{No aplica}.
\end{itemize}
\end{tencuenta}
\begin{algosalemal}
\begin{itemize}
  \item Si validaste por error, puedes devolver el documento con una observación mientras la etapa siga abierta.
  \item Si una etapa se cerró antes de tiempo, usa \boton{Reabrir etapa} en el expediente.
  \item Un trámite en pausa no se puede validar: aparece en la bandeja como \emph{Abandonado} hasta que el pasante lo retome.
\end{itemize}
\end{algosalemal}
